% =====================================================================
% JudgeSense - single-file build for Overleaf.
%
% Generated by flattening paper/main.tex and its \input files. The
% section boundaries are marked below; this file is now the source of
% truth for Overleaf, so edit it here.
%
% ANONYMITY: \blindtrue below produces the anonymous review copy.
% Comment that line out for the camera-ready. Do not upload a named PDF
% for review - TMLR review copies are anonymous.
%
% The results table is inlined from tables/main_results_v2.tex, which is
% GENERATED by scripts/regenerate_results.py from committed raw outputs.
% If the numbers change, regenerate there and re-flatten; do not hand-edit
% the table, or the paper and the data will drift apart.
% =====================================================================
% =====================================================================
% DO NOT UPLOAD THIS FILE TO OVERLEAF.
%
% This is the authored source. It \input's the generated results table
% from ../tables/, which does not exist inside an Overleaf project, so
% uploading it fails with "File `../tables/main_results_v2.tex' not
% found" and produces no PDF.
%
% Upload paper/overleaf_upload/main.tex instead -- same content, table
% inlined, no external inputs. Regenerate it after editing this file:
%     python scripts/build_overleaf_bundle.py
% =====================================================================
% ---------------------------------------------------------------------------
% JudgeSense --- TMLR submission.
%
% TMLR ships its own style file (tmlr.sty + tmlr.bst) from the author kit at
% https://jmlr.org/tmlr/. Drop those beside this file and swap the two marked
% lines below; the body needs no other change. Until then this builds with
% article so the document is checkable without the kit.
% ---------------------------------------------------------------------------
\documentclass[11pt]{article}          % TMLR: \documentclass{article} + \usepackage{tmlr}
\usepackage[margin=1in]{geometry}      % TMLR: delete (tmlr.sty sets the page)

\usepackage[utf8]{inputenc}
\usepackage[T1]{fontenc}
\usepackage{lmodern}                   % scalable fonts; microtype requires them
\usepackage{booktabs}
\usepackage{amsmath,amssymb}
\usepackage{graphicx}
\usepackage{microtype}
\usepackage[numbers,round]{natbib}
\usepackage[hidelinks]{hyperref}
\usepackage{url}

\title{JudgeSense: Measuring the Prompt Sensitivity of LLM-as-a-Judge}

% ANONYMITY SWITCH. TMLR review copies are anonymous, and a named review upload
% is a procedural rejection independent of content. This defaults to BLINDED so
% that the file cannot be submitted de-anonymised by forgetting a step; comment
% the line out (or build with \def\camera{}) for the camera-ready.
%
% Author order follows the superseded preprint (arXiv:2604.23478).
\newif\ifblind
\blindtrue                              % default: blinded, so a forgotten step
                                        % yields an anonymous file, never a
                                        % de-anonymised one
\ifdefined\camera\blindfalse\fi         % camera-ready: define \camera at build time

\ifblind
  \author{Anonymous authors\\Paper under double-blind review}
\else
  \author{
    Rohith Reddy Bellibatlu\thanks{Corresponding author.}\,$^{1}$ \quad
    Edward Raff$^{2,3}$ \quad
    Wenbin Zhang$^{1}$ \\[0.6em]
    {\normalsize $^{1}$Florida International University, Miami, FL, USA} \\
    {\normalsize $^{2}$CrowdStrike, New York, NY, USA} \\
    {\normalsize $^{3}$University of Maryland, Baltimore County, Baltimore, MD, USA}
  }
\fi
\date{}

\begin{document}
\maketitle

\begin{abstract}
Language models are routinely used to evaluate other language models, and an
evaluator earns that role by being reproducible. We ask a question the usual
validation cannot answer: does a judge return the same verdict when the same
request is worded differently? We introduce JudgeSense, a benchmark of 880 items
across four evaluation tasks drawn from human-labelled corpora, each issued under
two instruction templates that differ in wording and not in what they ask. The
metric never consults the ground-truth label --- a judge that is uniformly wrong
but uniformly consistent scores perfectly --- and it is reported against each
judge's own repeat-agreement ceiling, so ordinary decoding variance is separated
from sensitivity to wording rather than attributed to it. We release the
heuristic baselines that would beat each task if the construction were flawed,
with the scores they actually achieve, their intervals and their support counts,
alongside a threshold that a competent judge must clear for a task to count as
discriminating at all. Across three judges drawn from a single vendor family,
meaning-preserving rewording costs agreement in all ten of the twelve
judge--task cells that reached the declared support floor. The loss is largest on
the ordinal task for every judge and smallest on factuality, where the intervals
are consistent with an effect we declared to be practically meaningless --- a
null we report as a result. We also report a judge-specific
refusal behaviour that a naive harness would misreport as a format-adherence
failure, and the construction defects we found in our own benchmark, with the
measurements that exposed them.
\end{abstract}

% ===== 01_introduction =============================================
\section{Introduction}
\label{sec:intro}

Language models are now routinely used to evaluate the output of other language
models. The practice is cheap, fast, and correlates well enough with human
preference to have displaced human annotation across large parts of the
evaluation stack \citep{zheng2023judging}. Leaderboards are built on it, training
signals are derived from it, and papers report it as evidence.

An evaluator earns that role by being reproducible. If two researchers ask the
same judge the same question in different words and receive different verdicts,
the judge is not measuring the property it claims to measure. It is partly
measuring the phrasing. That failure is invisible in the usual validation, which
asks whether a judge agrees with humans on average. A judge can correlate well
with human ratings in aggregate while being unstable on any individual item, and
aggregate correlation cannot distinguish the two.

Instability under rephrasing is well documented for language models performing
ordinary tasks. Small, meaning-preserving changes to a prompt (formatting, ordering, the choice
between synonyms) move accuracy substantially
\citep{sclar2024quantifying, mizrahi2024state, elazar2021measuring}. The
evaluation literature has separately established that judges carry systematic
biases: they favour the first-presented candidate
\citep{wang2023fair, shi2024judging}, and length is a well-known confound --- in
the human preference votes we draw on, annotators chose the longer response in
roughly 70\% of usable comparisons. What has been measured less is the
property that matters most for a measuring instrument: whether the same judge,
asked the same question two ways, returns the same answer.

This paper measures that directly. We introduce JudgeSense, a benchmark of 880
items across four evaluation tasks, each issued under two instruction templates
that differ in wording and not in what they ask. The central quantity is not
accuracy. A judge that is wrong on every item but wrong \emph{identically} under
both phrasings is perfectly reproducible, and our metric assigns it a perfect
score. We are asking whether the instrument is stable, and treating whether it is
correct as a separate question.

\paragraph{Separating sensitivity from noise.}
A judge sampling at nonzero temperature disagrees with itself on the identical
prompt, so some part of any measured disagreement is ordinary decoding variance
rather than sensitivity to wording. Reporting raw agreement attributes all of it
to paraphrasing. We therefore issue each template a second time, unchanged, and
report the \emph{difference} between paraphrase agreement and repeat agreement.
The repeat arm is the judge's own reproducibility ceiling on that task, and the
gap below it is the quantity the paper is about. This distinction is not
cosmetic: in a pilot we discarded, a misconfigured decoding parameter produced a
repeat ceiling low enough to absorb the entire effect we report, and nothing in
the output distinguished it from a genuine one.

\paragraph{A benchmark has to be attackable before it is trustworthy.}
A measurement of judge stability is only as good as the items it is measured on.
If a task can be passed by a heuristic that ignores the intended construct, then
a high score says nothing about judgment. Always answer A. Pick the longer
candidate. Pick the passage sharing more words with the query. We therefore
run those heuristics against the released files and report what they score, with
intervals and support counts. None beats chance, and we report which control is
weakest and why. Equally, a benchmark on which \emph{nothing} works,
including a competent judge, is shortcut-free and useless; so we also declare a
threshold the best judge must clear for a task to count as discriminating at
all, and report each task against the rate a constant-answer judge achieves
rather than against uniform chance. Both directions are checked, because a construction that
suppresses exploitable signal can suppress genuine signal by the same mechanism.

\paragraph{This paper is a correction as well as a contribution.}
% The self-citation names all three authors in the bibliography, and the
% surrounding sentence claims the work as ours, so rendering it under review
% would deanonymise the submission as surely as the author block would.
An earlier version of this benchmark
\ifblind(reference withheld for review)\else\citep{bellibatlu2026judgesense}\fi{} was
withdrawn after review. Its items were
hardcoded in source while being described as benchmark-sourced; it shipped 75
unique items; it presented the correct candidate first in every pairwise item, so
a judge that always answered ``A'' was indistinguishable from a perfect one; and
one factuality template inverted the answer polarity, inflating measured
disagreement uniformly across every judge. We rebuilt from real corpora with
per-item provenance and fixed each defect. In doing so we found further defects of the same class in our own rebuild: a
template assignment that functioned as an answer key, candidate responses
carrying contradictory gold labels, and a length-balancing step that installed a
difficulty confound in place of the verbosity one it removed. We report those
too, with the measurements that exposed them. The failure mode we are documenting is not exotic. It is
ordinary, it is quiet, and it survives every check that does not specifically
look for it.

\paragraph{Contributions.}
\begin{enumerate}
  \item \textbf{A benchmark with attack surface measured, not asserted.} 880
    items from four human-labelled corpora with per-item provenance, released
    with the heuristic baselines that would beat it, their confidence intervals,
    and their support counts.
  \item \textbf{An endpoint measured against each judge's own noise ceiling.}
    $\Delta\mathrm{JSS}$, the gap between paraphrase agreement and a judge's own
    repeat-agreement ceiling. This demotes decoding variance rather than
    eliminating it, and we state the estimator's identity so a reader can see
    what it does and does not isolate. The endpoint, the estimator and the
    item-level clustering unit were pre-registered before the sweep; the support
    floors and the smallest effect of interest were added after the first
    judge's results and are reported as declared rather than pre-registered.
  \item \textbf{Refusal treated as a measurement rather than a nuisance.}
    Provider-flagged declines are recorded and reported separately from
    unparseable output, which are indistinguishable after parsing but are
    different facts about a judge.
  \item \textbf{Evidence that the ordinal task is the least stable.} Coherence
    carries the largest effect for all three judges and factuality the smallest,
    and on coherence the variation across judges is smaller than the variation
    across tasks within a judge. All three judges come from one vendor family, so
    this is a within-family observation and we do not claim it generalises.
\end{enumerate}
% ===== end 01_introduction =====
% ===== 02_related ==================================================
\section{Related work}
\label{sec:related}

\paragraph{LLM judges and their validation.}
Using a strong language model to score the output of another is now standard
practice, established at scale by MT-Bench and Chatbot Arena
\citep{zheng2023judging} and adopted across holistic evaluation efforts
\citep{liang2023holistic}. Validation of these judges is almost always
\emph{correlational}: a judge is trusted when its scores agree with human ratings
in aggregate \citep{chiang2023can, liu2023geval}. \citet{chiang2023can} find that
agreement is broadly positive but \emph{task-dependent}, which prefigures a result we
report below. That criterion is necessary and not sufficient. Aggregate
correlation is compatible with per-item instability, because errors that cancel
across a corpus still make any individual verdict unreliable, and a researcher who
reruns an evaluation with differently worded instructions has no guarantee of
reproducing it. Our measurement is orthogonal to accuracy by construction: it
never consults the ground-truth label, and a judge that is uniformly wrong scores
perfectly on it.

\paragraph{Known judge biases.}
A body of work documents specific, systematic ways judges fail. Position bias ---
preferring whichever candidate is presented first --- is strong enough to require
explicit correction \citep{wang2023fair}. \citet{shi2024judging} extend that
analysis across twelve judges and more than 100{,}000 evaluation instances,
finding position bias to be a stable, model-specific property \emph{modulated by
the quality gap between candidates}: it bites hardest when the two options are
close. That result directly motivates a construction choice we make in
Section~\ref{sec:benchmark}, where the preference split is matched on annotator
vote margin so that difficulty is not confounded with the property being measured.

Length is the other well-documented confound \citep{wang2023fair}, and rather
than lean on a citation for its magnitude we measure it in our own source data:
in the MT-Bench human preference votes we draw from, annotators preferred the
longer of the two responses in roughly 70\% of usable comparisons
(Section~\ref{sec:shortcuts}), so a judge that simply picked the longer answer
would score far above chance on an unfiltered sample.

These are findings about \emph{what} judges prefer. We take them as constraints
on benchmark construction rather than as our object of study: because position
and length bias are known, a benchmark that does not control them measures the
bias instead of the judge. We therefore present both candidate orderings for
every pairwise item and balance the length distribution exactly, then verify on
the released files that neither heuristic beats chance.

\paragraph{Prompt sensitivity in language models generally.}
That model behaviour moves under meaning-preserving reformulation is
well-established outside the evaluation setting. Formatting choices alone ---
delimiters, separators, whitespace --- swing accuracy by as much as 76 points on
some benchmarks \citep{sclar2024quantifying}; paraphrasing task instructions
changes measured capability enough to alter model rankings
\citep{mizrahi2024state}; and consistency under rephrasing has been treated as a
first-class property of pretrained models \citep{elazar2021measuring}. The
sensitivity has itself been assessed and benchmarked as a phenomenon
\citep{zhuo2024prosa, razavi2025benchmarking}, and \citet{razavi2025benchmarking}
report that current methods struggle to \emph{anticipate} which prompts will trip
a model up --- which is an argument for measuring sensitivity directly rather
than predicting it. \citet{wei2024systematic} study prompt-template effects on
LLM-judge reliability specifically, though restricted to alignment tasks and
under a different metric framework.

All of this concerns a model \emph{being tested}. Our contribution is to apply the
same lens to the model doing the testing. The distinction is not merely one of
subject: a judge does not generate an answer, it returns a label from a small
finite set, so its sensitivity is not about output fluency or content quality but
about decision stability across paraphrases. That is what makes a separate
measurement necessary rather than a re-use of existing ones.

\paragraph{Judging judges.}
Recent work evaluates LLM judges as objects in their own right.
\citet{thakur2024judging} assess thirteen judges on a clean alignment task and
document sensitivity to prompt complexity, length and other surface features
alongside a general tendency toward leniency; \citet{li2024llms} survey the area
broadly. These establish that judges are biased. What they do not provide is a
reportable, judge-level number that a practitioner can compute for a new judge
and act on.

The closest prior work to ours is \citet{arabzadeh2025human}, who study prompt
sensitivity in LLM-based relevance judgment directly, comparing human- and
LLM-authored prompt variants --- 90 prompts across binary, graded and pairwise
formulations, on three judge models --- and reporting Cohen's $\kappa$ against
TREC labels. The difference is what the measurement is referenced against. That
work evaluates labels produced under varying prompts against human relevance
gold, so it measures how rewording moves a judge's \emph{accuracy}; JSS never
consults the gold label and measures how rewording moves a judge's
\emph{agreement with itself}. The two are complementary and can diverge: a judge
whose accuracy is unchanged under rewording may still be flipping individual
verdicts in compensating directions, which our measurement sees and an
accuracy-referenced one does not. We note also that their relevance labels and
our relevance split both derive from TREC judgments, so the two are measuring
related material under different referents.

\paragraph{Why not simply evaluate under many prompts?}
The natural alternative to measuring sensitivity is to average it away.
\citet{mizrahi2024state} advocate multi-prompt evaluation as a methodological
default, and it is good advice for the evaluation being run. It is not a
substitute for a judge-level property. Multi-prompt evaluation must be re-run for
every new judge, over a prompt distribution chosen for that study; a fixed,
validated paraphrase set yields a number computed once and portable to any
evaluation context, which is what makes it usable as a selection criterion or a
regression test.

Variance-based characterisations \citep{sclar2024quantifying, zhuo2024prosa}
require gold labels, because they measure how accuracy moves across prompt
variants. Ours requires none, which is what allows it to be applied to a judge on
a corpus where no gold labels exist --- the situation a practitioner deploying a
judge is usually in. Ensemble judging \citep{li2024llms} manages sensitivity by
combining judges, but tells you nothing about whether any individual judge is
stable enough to use alone, which is the question a team choosing one judge
actually faces.

\paragraph{Sources and statistical machinery.}
The items come from existing human-labelled corpora: TruthfulQA for factual
accuracy \citep{lin2022truthfulqa}, SummEval for expert coherence ratings
\citep{fabbri2021summeval}, BEIR's TREC-COVID split for graded relevance
judgements \citep{thakur2021beir}, and MT-Bench's human preference votes
\citep{zheng2023judging}. No label in this benchmark is assigned or inferred by
us. Agreement is chance-corrected with Cohen's $\kappa$
\citep{cohen1960coefficient}, quadratic-weighted on the ordinal task, and all
intervals are cluster bootstrap \citep{efron1979bootstrap} resampled at the item,
since the two candidate orderings of a pairwise item and its repeat arms are not
independent observations. The judges evaluated are drawn from current
instruction-tuned and reasoning-tuned families \citep{openai2023gpt4,
gemini2023gemini, grattafiori2024llama3, qwen2024qwen25, jiang2023mistral,
deepseek2025r1}.
% ===== end 02_related =====
% ===== 03_benchmark ================================================
% ---------------------------------------------------------------------------
% Benchmark construction. Number-independent: everything here is a property of
% the artifact, not of any judge's behaviour, so it is stable across runs.
% ---------------------------------------------------------------------------
\section{The JudgeSense benchmark}
\label{sec:benchmark}

JudgeSense measures whether an LLM judge returns the same verdict when the same
evaluation request is worded differently. The construct is deliberately narrow:
we are not asking whether a judge is \emph{correct}, but whether it is
\emph{stable}. A judge that is consistently wrong is reproducible; a judge whose
verdict depends on which of two equivalent phrasings it received is not usable as
a measuring instrument, however accurate it looks on average.

\subsection{Tasks and sources}
\label{sec:sources}

The benchmark covers four evaluation tasks, each drawn from an existing
human-labelled corpus. No item is authored by us, and no label is assigned or
inferred by us; every label is the one the source corpus already carried.

\begin{description}
  \item[Factuality] (250 items, TruthfulQA). A statement is judged accurate or
    inaccurate. Items whose question is ill-posed, and whose only defensible
    answer is a refusal to answer, are excluded at load time.
  \item[Coherence] (250 items, SummEval). A summary is rated 1--5 for
    coherence against expert annotations. Summaries with identical text but
    conflicting expert labels are dropped rather than merged.
  \item[Relevance] (250 items, 500 rows, BEIR TREC-COVID). Given a topic, the
    judge picks the more relevant of two passages. \emph{Both} candidates carry
    a human relevance grade: the positive is graded fully relevant (2) and the
    distractor is one a human assessor explicitly rejected (0). Partially
    relevant documents (grade 1) are used for neither role.
  \item[Preference] (130 items, 260 rows, MT-Bench human judgments). The judge
    picks the response humans preferred. Only comparisons with at least two
    annotator votes and a clear majority are used.
\end{description}

The relevance construction is the substantive departure from the usual practice
of drawing distractors from a positives-only collection. There, a negative can
only be a document \emph{absent} from an incomplete relevance judgment set, so a
topically similar distractor may in fact be relevant but unjudged, and the item
penalises a defensible answer. TREC-COVID carries graded judgments including
tens of thousands of explicit non-relevant assessments, so here both candidates
are human-labelled and the distractor is one a human affirmatively rejected.

Every shipped item records its source dataset, split, per-record identifier, and
the source configuration where the source defines one, so any item can be traced
to the row it came from.
Loaders raise rather than substitute when a source is unavailable; there is no
fallback path that could silently generate an item.

\subsection{Paraphrase design}

Each item is issued under two of five instruction templates, intended to differ
in wording and not in what they ask. That intent is a construct assumption, and
if it fails the benchmark stops measuring what it claims to: were one template
to ask a subtly different question, a judge answering both correctly would be
scored as unstable, and semantic drift rather than prompt sensitivity would
explain the result. We therefore state plainly what we have checked and what we
have not.

We audit the released files themselves, across all four tasks, recovering each
template from the shipped prompts and the payload recorded alongside them rather
than trusting the generator. The audit is offline and deterministic;
it makes no model calls, and it is run in continuous integration. Four
properties are enforced. \emph{Label space}: the set of admissible answers each
template requests, parsed from the instruction with the item text and the
candidate frame removed, is identical within every pair (\{YES, NO\} for
factuality, $\{1,\dots,5\}$ for coherence, \{A, B\} for both pairwise tasks).
\emph{Polarity}: no answer token is bound to opposite conclusions across
templates, no template asks the inverted question, and where a template glosses
its scale endpoints the low anchor is the negative one. \emph{Construct}: every
template names a term from its own task's declared vocabulary and none from
another task's. \emph{Non-triviality}: no pair is a near-duplicate, so every
pair is a real manipulation. All four hold on the shipped templates for all four
tasks. This supersedes our earlier check, which iterated only the factuality and
coherence templates, leaving the two pairwise tasks entirely uncovered, and
tested substring presence rather than the answer set requested.

Two further properties are measured and reported rather than enforced, because
we cannot show they are harmless. Templates place the payload at different
points in the instruction: in relevance, T1 opens the query at $46.4\%$ of the
template against T3 at $9.0\%$, a within-pair divergence of up to $37.4$
percentage points. Instruction reading difficulty also varies, by up to $3.7$
Flesch--Kincaid grade levels within a pair (preference T3--T4). Either could
make one arm intrinsically harder, and neither is controlled by our design.

We do not claim these checks establish semantic equivalence. They bound specific
failure modes; nothing computable from the template strings settles whether a
competent reader maps two wordings onto the same question. Three residuals
remain open. The audit groups surface forms into one construct by an authored mapping.
Factuality templates ask variously about a statement being ``factually
correct'', ``accurate'' or worth ``fact-checking''; preference templates ask
which response is ``better'', ``of higher quality'' or ``stronger''. Whether
those are synonyms for a given judge is an assumption, not a measurement. Implicit thresholds may shift between such
phrasings even when the construct does not. And tokenisation or pretraining
frequency may favour one phrasing for a particular model. Semantic drift is
therefore reduced as an alternative explanation for measured sensitivity, not
eliminated; establishing equivalence would require human annotation or a
model-in-the-loop study we have not run.

The polarity constraint is a correction of our own earlier design. An earlier version of
this benchmark included a factuality template that inverted the answer polarity
(``does this response contain factual errors? answer NO if accurate''). Under a
single shared label space, ``YES'' then denoted opposite conclusions in different
arms, so a judge that answered both arms correctly was recorded as having
disagreed with itself. That is a label-convention defect in the benchmark rather
than a property of any model, and it measures a different construct: handling a
polarity flip is a comprehension task, whereas we ask whether equivalent
phrasings yield equivalent verdicts.

We state the direction of the resulting bias rather than leave it implicit.
Polarity-inverted pairings are among the hardest items in the set, so excluding
them can only push measured consistency \emph{upward}; our stability figures are
conditioned on that exclusion and should be read accordingly.

The alternative --- retaining inverted templates and remapping each template's
raw answers onto one canonical decision space before comparison --- is
implemented in the released code, but we cannot report it here and we say why
rather than deferring it. The v2 construction removes polarity inversion at
source: all twenty shipped templates request the same answer space in the same
direction, verified by the audit above. There is consequently nothing in this
dataset for the remap to operate on, and no remapped analysis is possible without
building a polarity-inclusive split we have not built. The machinery is released
for anyone who wants to construct one; the quantification of what the exclusion
costs is future work, not a result we are withholding.

The earlier version did have that quantification, and its shape is worth
recording even though its magnitude rested on data we have withdrawn. Including
polarity-inverted pairings produced an apparent flip rate that was near-identical
across every judge tested. That uniformity is the diagnostic: a rate that does
not vary with the model is a property of the instrument rather than of the
models, and it is what identified the template as a label-convention defect
rather than a finding about judges. We keep the inference and discard the number.

A limitation of the assignment follows from how templates are allocated. Each
item is issued under exactly one of the ten unordered template pairs, balanced
across pairs: 25 items per pair for factuality, coherence and relevance, and 13
for preference. Every item is therefore seen under exactly two of the five
templates, and no item is seen under more than one pair. Template identity is
consequently nested perfectly inside item identity: the items behind template
$T_i$ and those behind $T_j$ are disjoint sets, so any apparent difference
between two templates is equally a difference between the items they were shown,
and no estimator applied to this design can separate the two. Statements of the
form ``template A is more stable than template B'' are not identifiable here,
and we make none. What the design does support is the within-item contrast it
was built for: the two arms of a pair differ only in wording, with the item held
fixed, so pair-level sensitivity is unaffected by this confound. Separating
per-template effects would require crossing items over multiple template pairs,
which this release does not do.

\subsection{Position control}

Both pairwise tasks present every item twice, once in each candidate ordering.
The correct answer therefore sits at position A in exactly half the rows, and a
judge that ignores content and always answers ``A'' scores $0.500$ rather than
$1.000$.

This too corrects an earlier version, in which the correct candidate always
appeared first. Under that construction a position-anchored judge was
indistinguishable from a perfect one, and most judges scored a degenerate
$\mathrm{JSS}=1.000$ that reflected the layout rather than any judgment.

\subsection{Shortcut controls}
\label{sec:shortcuts}

A benchmark measures what it claims to measure only if it cannot be passed by a
heuristic that ignores the intended construct. We therefore evaluate the shipped
splits against judges that use no judgment at all. Each is run over the released
files, not over the loaders, so these are properties of the artifact:

\begin{itemize}
  \item \textbf{Position only} (always answer A): $0.500$ on relevance
    ($n=500$) and $0.500$ on preference ($n=260$), by construction of the swap
    design.
  \item \textbf{Length only} (pick the longer candidate): $0.482$ on relevance
    ($n=494$), $0.508$ on preference ($n=256$).
  \item \textbf{Lexical overlap only} (pick the candidate sharing more terms with
    the query): $0.540$ on relevance ($n=252$).
\end{itemize}

Support counts fall below the row counts because a heuristic that ties has no
signal to act on; tied rows are excluded rather than split, which is why the
overlap control rests on roughly half the relevance rows. All three controls are
computed from the released files by \texttt{scripts/shortcut\_controls.py}, whose
output is committed, so every figure above is reproducible without rerunning any
judge.

None of these exceeds chance in either direction, though the overlap control is
the weakest of the three and we do not overstate it: its 95\% Wilson interval is
$[0.478, 0.600]$, so the data are consistent with a residual lexical signal worth
about $0.57$. It also reads only half the relevance rows, because the remainder
tie on the overlap statistic and a heuristic with no signal on a row is not
evidence about that row. The length and position controls are far tighter.

The length control is enforced rather than observed. Human annotators in the
preference corpus preferred the longer response in roughly 70\% of usable
comparisons, so an unfiltered sample rewards verbosity directly; we hold the
winner-longer share at exactly 50\% by taking an equal number of items from each
length bucket. That $50\%$ is a construction invariant of the item set and is a
different quantity from the $0.508$ above, which is what a length-following judge
scores on the rendered rows once ties are dropped.

The preference split ships 130 items rather than the round 250, and we state the
full attribution rather than the net gap. Of the source comparisons, 607 failed
the decisive-vote label rule or carried no strict majority, one was duplicate
content, 67 were dropped because their gold contradicted another item's gold on
shared candidate text, and 102 more were dropped to hold the length balance
exactly. Each filter removes a way the benchmark could be passed or contradicted
without judging, and the last two were added after we found both failures in our
own shipped data. We regard the reduced split as the correct trade: the
alternative is a larger set that a verbosity heuristic beats, and one that scores
a consistent judge as wrong whenever the same response appears on both sides of
two different comparisons.

Balancing on length alone would install a second confound in place of the first.
Because the winner-shorter bucket ships whole and only the winner-longer bucket
is cut, any criterion used to make that cut applies to one bucket and not the
other. Cutting by vote margin --- keeping the most contested comparisons, as an
earlier build did --- left the buckets matched on length but far apart on
annotator agreement: mean vote-margin ratio $0.567$ against $0.745$, and $16\%$
against $54\%$ unanimous. A judge that is stronger on close calls would then
read as length-biased. We therefore \emph{stratify}: the winner-shorter bucket
fixes the target distribution of vote shapes (absolute margin and total votes),
and the winner-longer bucket is filled stratum by stratum to match it. The
shipped buckets agree to $0.0003$ in mean vote-margin ratio ($0.8603$ against
$0.8605$), have identical medians, and contain the same number of unanimous
comparisons ($49$ each).

Matching stops the split from collapsing onto coin-flips, but the opposite
failure is equally real and less obvious. The filter that removes contradictory
gold is a greedy selection over a conflict graph, so the order in which items are
considered determines both which survive and how hard the survivors are. Keeping
the firmest human majorities first is the intuitive rule and is worse on both
counts, measured on this pool: $106$ items at $90.6\%$ unanimous, against $130$
items at $75.4\%$ when the most contested are kept first. A split on which the
annotators were nearly always unanimous cannot separate a competent judge from a
mediocre one --- both score near the top --- which is the ceiling effect that
made the earlier version of this benchmark uninformative. We therefore resolve
conflicts contested-first, and report the resulting distribution rather than only
the balance: $75.4\%$ of shipped comparisons carry a unanimous human majority and
the rest are genuine close calls.

The lexical control required a similar correction. Selecting each distractor to
maximise its term overlap with the topic --- the obvious way to make an item
``hard'' --- made distractors \emph{denser} in query terms than the genuinely
relevant passages, so the heuristic ``pick the passage with \emph{lower} overlap''
recovered the answer far above chance. Matching the distractor's retrieval score
to the positive's, rather than maximising it, removes overlap as a signal in
either direction. We report the control two-sided for this reason: a split whose
answer is reliably the \emph{shorter} or \emph{less overlapping} candidate is
exploited exactly as cheaply as one biased the other way.

\subsection{Release and integrity}

The dataset is released with Croissant metadata, per-item provenance, and a
content hash for each split that excludes retrieval timestamps, so a rebuild can
be verified as identical in content without being byte-identical in metadata. A
32-check audit gate (eight checks over each of four splits) runs in continuous integration and fails the build on
duplicate content, unbacked provenance claims, label degeneracy, insufficient
clustering support, contradictory ground truth, or a length shortcut exceeding
its bound. Checks that cannot read the data they are meant to inspect fail
rather than pass, so a check that silently measures nothing cannot be mistaken
for a check that found nothing wrong.
% ===== end 03_benchmark =====
% ===== 04_metrics ==================================================
% ---------------------------------------------------------------------------
% Metric definitions and the analysis plan. Number-independent.
% ---------------------------------------------------------------------------
\section{Measuring sensitivity}
\label{sec:metrics}

\subsection{Judge Sensitivity Score}

For an item issued under two equivalent templates, let $d_1(p)$ and $d_2(p)$ be
the judge's parsed decisions on pair $p$. The Judge Sensitivity Score is the
agreement rate

\[
  \mathrm{JSS} = \frac{1}{|P|}\sum_{p \in P} \mathbb{1}[d_1(p) = d_2(p)],
\]

over the set of scored pairs $P$. JSS never consults the ground-truth label. A
judge that is wrong on every item but wrong \emph{identically} under both
phrasings scores $1.0$, and this is intended: it is a measure of reproducibility,
not of accuracy. We report accuracy separately, and the two answer different
questions. On the pairwise tasks the companion measure is \emph{position-corrected}
accuracy: whether the judge selected the side carrying the ground truth, given
that the swap design places it at position A in exactly half the rows. On the
pointwise tasks it is ordinary accuracy against the source label. The
distinction matters: a judge that always answers ``A'' has an ordinary accuracy
that depends only on the layout, while its position-corrected accuracy is
$0.500$ and its behaviour is immediately visible.

Raw agreement is inflated when a judge's output distribution is compressed, since
a judge answering ``YES'' almost always agrees with itself. We therefore report
the chance-corrected form (Cohen's $\kappa$; quadratic-weighted for the
ordinal coherence task) alongside every raw figure.

We claim no novelty for this quantity and say so explicitly. Agreement rate under
paraphrase is the same measurement \citet{elazar2021measuring} use for consistency
in pretrained models, that \citet{arabzadeh2025human} compute as Cohen's $\kappa$
against relevance labels, and that \citet{zhuo2024prosa} express as a variance;
JSS is mathematically of that family. What is new here is the benchmark, the
systematic application to the model doing the judging rather than the model being
judged, and the endpoint defined next --- which references that agreement against
the judge's own repeat ceiling, so a figure that would otherwise mix wording
sensitivity with decoding noise separates the two.

\subsection{The repeat baseline, and why it is the endpoint}
\label{sec:repeat}

Raw JSS below $1.0$ does not by itself demonstrate sensitivity to \emph{wording}.
A judge sampling at nonzero temperature, or one served through a
non-deterministic backend, disagrees with itself on the identical prompt. Some
portion of any measured disagreement is therefore ordinary decoding noise, and a
paper that reports raw JSS attributes all of it to paraphrasing.

We separate the two by issuing one template a second time, unchanged. The
resulting $\mathrm{JSS}_{\text{rep}}$ is that judge's reproducibility ceiling on
that task: the agreement it achieves when nothing varies. The primary endpoint is
the difference

\[
  \Delta\mathrm{JSS} = \mathrm{JSS}_{\text{paraphrase}} - \mathrm{JSS}_{\text{rep}},
\]

which is negative when rewording costs more stability than resampling does. This
is the quantity the paper's claim is about, and it was fixed in a pre-registered
analysis plan committed to the repository (commit \texttt{b1b507a}) before the
full sweep was run.

We state the history precisely, because a pre-registration is worth only what is
true of it, and ours is partly post-hoc.

At commit \texttt{b1b507a} the plan fixed the endpoint, the estimator, the
item-level clustering unit, and one decision rule: an effect is present when the
interval excludes zero. Those four are genuinely pre-committed, and they preceded
the sweep --- though not the first API call, since a five-row-per-cell smoke wave
had already been run against three judges to exercise the harness. That wave
produced no reportable estimate and none of its output enters the results.

Three further elements were added at commit \texttt{82ea522}, roughly seventy
minutes after the first full judge's results were committed: the support floors,
the multiplicity plan, and the discrimination thresholds. They are \emph{not}
pre-registered and we do not describe them as such. Two of the four
discrimination thresholds were set while the corresponding accuracies were
already visible in the repository, which we regard as disqualifying them as
confirmatory tests; they are reported as post-hoc descriptive checks. The other
two, on the pointwise tasks, were set before any pointwise accuracy had been
computed.

One correction is recorded in the plan itself: the direction of subtraction in
the endpoint's prose was stated backwards relative to the estimator, which had
been implemented and tested beforehand. The estimator was unchanged and both
orderings were visible in the committed output throughout.

The multiplicity plan added at \texttt{82ea522} is not applied in this paper. It
called for a pooled contrast per task with Holm correction, per-cell intervals
demoted to exploratory under Benjamini--Hochberg control, and a reported minimum
detectable effect. We report per-cell intervals without adjustment and label
them accordingly: the twelve cells are descriptive, not confirmatory, and no
judge is ranked against another on the strength of an unadjusted cell.

\subsection{Uncertainty}

All intervals are 95\% cluster bootstrap intervals over 2{,}000 resamples, with
the item as the resampling unit. Item-level clustering is required rather than
preferred: a pairwise item contributes two rows (one per candidate ordering) and
its repeat arm nests inside it, so resampling rows independently treats
correlated observations as independent and understates uncertainty. The
clustering unit is a mandatory argument in our implementation, with no default,
so it cannot be omitted by accident.

An effect is reported as present for a judge--task cell when the interval for
$\Delta\mathrm{JSS}$ excludes zero. The direction is read from the sign.

\subsection{Outcome categories: verdict, refusal, malformed}
\label{sec:refusal}

Each call terminates in exactly one of three outcomes.

A \textbf{verdict} is a response the strict parser maps to a label in the task's
answer space. A \textbf{refusal} is a response the provider itself flags as
declined, typically with empty content. Providers spell this differently:
\texttt{stop\_reason="refusal"} for Anthropic, \texttt{finish\_reason=
"content\_filter"} for OpenAI-compatible APIs, \texttt{SAFETY} or
\texttt{RECITATION} for Google. The accepted set is therefore enumerated in a
single place in the released code and matched case-insensitively, and the raw string the
provider returned is retained on every record. Recognising only one family's
spelling would report a refusal rate of zero for every other family regardless of
what those judges did, and would silently charge their declines as paraphrase
disagreement.
\textbf{Malformed output} is a completed response the parser cannot map. Parsing
is strict and never coerces: an unmappable response is recorded with its raw text
retained, never rounded to the nearest label.

Malformed output and refusal are treated differently, and the asymmetry is
deliberate. A malformed arm is \emph{charged as disagreement} and stays in the
support: the model completed a response, so the pair reflects a judgment the
parser could not read rather than one that was never rendered. A refused arm is
removed from the support, because no judgment was rendered at all. Malformed-output
rate is reported per cell alongside the refusal statistics below, so neither is
absorbed into the other.

Refusals and malformed output are indistinguishable after parsing --- both yield
no usable label --- but they are different measurements, and we do not merge
them. A refusal occurs \emph{upstream} of any judgment: the provider halted
before the model rendered a verdict. Scoring it as paraphrase disagreement
asserts that the judge produced two conflicting judgments, which it did not.
Treating refusal as a third label is worse still, since a judge that refuses
everything would then score $\mathrm{JSS} = 1.0$.

We therefore compute JSS over pairs in which both arms returned verdicts, and
report the support size $n$ alongside every figure. Refusal is reported as its
own construct, per cell:

\begin{itemize}
  \item \emph{refusal rate}: refused calls over all calls;
  \item \emph{refusal discordance rate} (RDR): pairs in which exactly one arm was
    refused;
  \item \emph{consistent refusal rate}: pairs in which both arms were refused.
\end{itemize}

RDR is itself a sensitivity statistic, and the most interesting of the three: a
nonzero RDR means a meaning-preserving rewording changed whether the judge was
willing to judge at all. That is prompt sensitivity of the willingness-to-answer
layer, and it is invisible to any analysis that scores refusals as disagreement
or discards them.

Conditioning on verdict pairs restricts a refusing judge to the items it agreed
to judge, which raises the fair objection that judges are then compared on
different subsets. We answer it with a measurement rather than an argument: every
cross-judge comparison in a task with nonzero refusals is recomputed on the
common support --- the items every compared judge answered under both phrasings
--- and we additionally report a refusal-inclusive figure in which every refused
arm counts as disagreement, the most punitive available reading. Both appear
beside the conditioned figure, so the effect of the conditioning is visible
rather than asserted.

\subsection{Decoding budget}
\label{sec:budget}

Every judge is run with an identical output token limit. This is not a neutral
default: a class-dependent limit --- a tight cap for instruction-tuned judges and
a generous one for reasoning-tuned judges, on the reasoning that the latter emit
chain-of-thought --- confounds judge class with inference configuration, so any
difference between the classes admits an explanation that has nothing to do with
the models. It also truncates real responses, since a judge that emits any
preamble before its verdict exhausts a tight cap mid-sentence and the fragment is
retained as its answer. Disagreement so produced is a property of the cap rather
than of the judge.

We deliberately quote no rate for this. Every committed run was made under the
matched policy, so we hold no measurement of the cap-termination rate under the
class-dependent one, and stating a rate we have not measured is the failure this
paper exists to avoid repeating. Until the ablation below is run, the case for
the matched budget rests on the confound being structural, not on its size.

The matched budget removes the explanation by construction. Because a token
limit is a ceiling and not a target, compliant judges that emit a single-token
answer are unaffected in cost; only verbose non-compliers cost more, and those
are precisely the calls that were previously being truncated into noise. We
have not quantified the budget's effect. Doing so would require a subsample run
under the class-dependent configuration, reporting the difference in
$\Delta\mathrm{JSS}$ and the cap-termination rate under each policy with their
denominators. We did not run it, and we make no claim about the effect's size.
% ===== end 04_metrics =====
% ===== 05_results ==================================================
\section{Results}
\label{sec:results}

We report three judges from one vendor family (Claude Haiku 4.5, Sonnet 4.5 and
Opus 4.7) on the four tasks. Factuality, coherence and relevance are
complete for all three; the preference split is complete for Haiku only, and
Section~\ref{sec:results-limits} states what that costs. Every figure below is
produced by \texttt{scripts/regenerate\_results.py} from the committed raw
outputs; none is transcribed.

% Generated by scripts/regenerate_results.py from the committed raw outputs.
% Never edit by hand: re-run the script instead, or the paper and the data drift.
% ---- begin main_results_v2.tex (inlined; edit here, this is the single source) ----
\begin{table}[t]
\centering
\caption{Paraphrase sensitivity per judge and task. $\Delta\mathrm{JSS} = \mathrm{JSS}_{\text{para}} - \mathrm{JSS}_{\text{rep}}$ is the pre-registered endpoint; a negative value means rewording costs more agreement than re-issuing the identical prompt. Intervals are 95\% item-clustered bootstrap over 2{,}000 resamples. Cells below the declared support floor of 100 clusters report no endpoint rather than one computed from too few. $r$ is the provider-reported refusal rate and $m$ the malformed-output rate over the arms the judge attempted, so the two do not overlap. The JSS column is computed over all rows of the cell; $\mathrm{JSS}_{\text{rep}}$ and $\Delta\mathrm{JSS}$ are computed over the canonical ordering only, so the two columns do not subtract to the printed delta where a cell carries refusals or a swapped-ordering imbalance.}
\label{tab:main}
\begin{tabular}{llrrrlrr}
\toprule
Judge & Task & $n$ & JSS & $\mathrm{JSS}_{\text{rep}}$ & $\Delta\mathrm{JSS}$ [95\% CI] & $r$ & $m$ \\
\midrule
claude-haiku & factuality & 250 & 0.956 & 1.000 & $-0.044$ [-0.072, -0.020] & 0.000 & 0.000 \\
 & coherence & 250 & 0.792 & 0.998 & $-0.206$ [-0.258, -0.154] & 0.000 & 0.000 \\
 & relevance & 500 & 0.930 & 0.990 & $-0.058$ [-0.086, -0.032] & 0.000 & 0.010 \\
 & preference & 260 & 0.896 & 0.965 & $-0.081$ [-0.127, -0.038] & 0.000 & 0.031 \\
\addlinespace
claude-opus-4-7 & factuality & 250 & 0.960 & 0.990 & $-0.030$ [-0.056, -0.006] & 0.000 & 0.000 \\
 & coherence & 250 & 0.736 & 0.916 & $-0.180$ [-0.236, -0.128] & 0.000 & 0.000 \\
 & relevance & 500 & 0.964 & 0.998 & $-0.036$ [-0.063, -0.013] & 0.051 & 0.000 \\
 & preference & 170 & 0.963 & --- & \emph{support below floor} & 0.000 & 0.056 \\
\addlinespace
claude-sonnet & factuality & 250 & 0.952 & 0.990 & $-0.038$ [-0.062, -0.018] & 0.000 & 0.008 \\
 & coherence & 250 & 0.828 & 0.990 & $-0.162$ [-0.210, -0.118] & 0.000 & 0.000 \\
 & relevance & 500 & 0.584 & 0.789 & $-0.158$ [-0.206, -0.110] & 0.247 & 0.224 \\
 & preference & 44 & 0.559 & --- & \emph{support below floor} & 0.000 & 0.432 \\
\bottomrule
\end{tabular}
\end{table}
% ---- end main_results_v2.tex ----

\subsection{Paraphrasing costs agreement beyond decoding noise}

Table~\ref{tab:main} reports $\Delta\mathrm{JSS} = \mathrm{JSS}_{\text{para}} -
\mathrm{JSS}_{\text{rep}}$. In all ten cells with sufficient support the interval
excludes zero and the effect is negative: a meaning-preserving rewording costs
more agreement than re-issuing the identical prompt does. One of those ten, Sonnet on relevance, we judge uninterpretable as evidence about
wording, for
reasons given in Section~\ref{sec:results-sonnet-relevance}. We keep it in the
count and in the table rather than deleting it, and exclude it from the
comparative claims that follow, so a reader can see exactly which conclusions
rest on it and which do not.

We state immediately what that uniformity does and does not establish, because
the estimator makes a weaker claim than a ten-for-ten result suggests. Under the
strict policy, $\Delta\mathrm{JSS}$ estimates
$-\tfrac{1}{2}\,\mathbb{E}_{\text{item}}\big[\sum_k (p_k - q_k)^2\big]$, where
$p$ and $q$ are the judge's response distributions under the two templates. It is
non-positive for any two prompts that do not induce identical distributions, so
a point null of exactly zero is false before any data are collected and rejecting
it is a statement about power. What the sign test establishes is only direction.
The quantity that carries the argument is the magnitude, which we assess against
the smallest effect declared of interest rather than against zero
(Section~\ref{sec:results-sesoi}).

The repeat ceiling ranges from $0.789$ to $1.000$. Nine of the ten cells sit at
$0.916$ or above, which is what makes the paraphrase gap interpretable for them:
at a matched budget and temperature zero these judges are close to deterministic
on a fixed prompt.

The tenth is treated separately in
Section~\ref{sec:results-sonnet-relevance}, because a ceiling of $0.789$ means
the judge disagrees with itself on more than a fifth of byte-identical prompts.
That is not a small caveat. In a pilot we discarded, a misconfigured decoding
parameter produced a ceiling of $0.864$, higher than this one, and we discarded
it precisely because a ceiling that low can absorb the effect being measured. We
apply the same standard to our own shipped cell.

A high ceiling carries an implication we would rather state than leave a reader
to derive. In seven of the ten cells it sits at $0.990$ or above, so the
subtraction defining our endpoint is nearly a no-op there and
$\Delta\mathrm{JSS} \approx \mathrm{JSS} - 1$. The repeat arm does substantial
work in only three cells: Opus on coherence ($0.916$), Haiku on preference
($0.965$), and the flagged Sonnet relevance cell ($0.789$). We collect it
everywhere regardless, because which cells those will be is not knowable in
advance, and because a judge run at a production temperature above zero has a
lower ceiling throughout. The correction matters more in deployment than our own
numbers make it look.

\subsection{Effect magnitudes against the declared threshold}
\label{sec:results-sesoi}

We declared $|\Delta\mathrm{JSS}| < 0.02$ to be not practically meaningful. That
threshold was set after the first judge's results were in hand (see
Section~\ref{sec:metrics}), so it is a declared standard rather than a
pre-registered one; we state it because a reader should know which of our
thresholds preceded the data and which did not. Applying that to intervals rather than to point estimates, as an
equivalence-style reading requires, separates the cells sharply:

\begin{center}
\begin{tabular}{llrl}
\toprule
Judge & Task & $\Delta\mathrm{JSS}$ & Interval vs.\ the $0.02$ threshold \\
\midrule
Haiku  & coherence  & $-0.206$ & entirely beyond \\
Opus   & coherence  & $-0.180$ & entirely beyond \\
Sonnet & coherence  & $-0.162$ & entirely beyond \\
Sonnet & relevance  & $-0.158$\,$^{\dagger}$ & entirely beyond \\
Haiku  & preference & $-0.081$ & entirely beyond \\
Haiku  & relevance  & $-0.058$ & entirely beyond \\
\midrule
Haiku  & factuality & $-0.044$ & upper edge at the threshold \\
Sonnet & factuality & $-0.038$ & upper edge inside \\
Opus   & relevance  & $-0.036$ & upper edge inside \\
Opus   & factuality & $-0.030$ & upper edge inside \\
\bottomrule
\end{tabular}
\end{center}

\noindent{\footnotesize $^{\dagger}$ Clears the threshold arithmetically but is
not interpretable as evidence about wording
(Section~\ref{sec:results-sonnet-relevance}).}

All ten supported cells appear here. Six clear the threshold unambiguously,
though one of those six is the flagged cell and we do not count it as support.
Four are consistent with an effect we declared to be practically meaningless: on
factuality the data do not distinguish a real cost from a negligible one for any
judge. We report factuality as a null on that reading. Our pre-registration does
commit us to treating a null as a result, and this is one.

\subsection{Sensitivity is larger on the ordinal task than on the others}

Coherence carries the largest effect for every judge, and factuality the
smallest. Between those two the picture is less clean than a compact ordering
suggests, and we do not claim more than the design supports:

\begin{center}
\begin{tabular}{lcccc}
\toprule
Judge & coherence & preference & relevance & factuality \\
\midrule
Haiku 4.5   & $-0.206$ & $-0.081$          & $-0.058$          & $-0.044$ \\
Sonnet 4.5  & $-0.162$ & \emph{no support} & $-0.158$\,$^{\dagger}$ & $-0.038$ \\
Opus 4.7    & $-0.180$ & \emph{no support} & $-0.036$          & $-0.030$ \\
\bottomrule
\end{tabular}
\end{center}

\noindent{\footnotesize $^{\dagger}$ Not interpretable as sensitivity to wording;
see Section~\ref{sec:results-sonnet-relevance}.}

Only Haiku has all four tasks, so a complete four-task ordering exists for one
judge. For the other two the available tasks are consistent with coherence
first and factuality last, but the adjacent pairs are not separated: Sonnet's
coherence and relevance differ by $0.004$ with almost coincident intervals, and
Opus's relevance and factuality differ by $0.006$ with overlapping intervals.

The comparison of spreads is likewise more equivocal than one pairing suggests,
so we report all of them:

\begin{center}
\begin{tabular}{lrlr}
\toprule
Within task, across judges & spread & Within judge, across tasks & spread \\
\midrule
factuality & $0.014$ & Sonnet (3 tasks) & $0.124$ \\
coherence  & $0.044$ & Opus (3 tasks)   & $0.150$ \\
relevance  & $0.122$ & Haiku (4 tasks)  & $0.162$ \\
\bottomrule
\end{tabular}
\end{center}

Within-judge spread exceeds within-task spread in every pairing, but the margin
is decisive only for factuality and coherence. Relevance ($0.122$) is
indistinguishable from Sonnet's within-judge spread ($0.124$), and both figures
depend on the same flagged cell. The defensible claim is narrower than a general
one: \emph{on the ordinal task, sensitivity is large and consistent across three
judges, and larger than the variation between those judges}. We do not claim
that task dominates judge in general, and a design with judges from more than one
vendor would be needed to support it.

Coherence being the most affected task is consistent with its structure. It is
the only ordinal task, so a judge has five admissible answers rather than two,
and a rewording that shifts a judgement by one scale point counts as
disagreement. The quadratic weighted kappa, which charges an adjacent-category
disagreement less than a distant one, scores $0.859$--$0.874$ against an
unweighted kappa of $0.594$--$0.727$: most coherence disagreements are one-point
moves.

\subsection{The tasks discriminate, and coherence only narrowly}
\label{sec:results-discriminate}

A benchmark on which no heuristic beats chance is worthless if no competent judge
beats chance either. Section~\ref{sec:shortcuts} established the floor. The
complementary check requires the best judge on each task to exceed a threshold;
we report it against the \emph{majority-class} rate rather than uniform chance,
because a constant-answer judge achieves the former and quoting the latter
overstates the margin.

\begin{center}
\begin{tabular}{lrrrr}
\toprule
Task & threshold & best observed & majority class & margin \\
\midrule
factuality & $0.75$ & $0.944$ (Opus)   & $0.500$ & $+0.444$ \\
relevance  & $0.70$ & $0.882$ (Opus)   & $0.500$ & $+0.382$ \\
preference & $0.65$ & $0.873$ (Haiku)  & $0.500$ & $+0.373$ \\
coherence  & $0.40$ & $0.452$ (Sonnet) & $0.348$ & $+0.104$ \\
\bottomrule
\end{tabular}
\end{center}

Three tasks discriminate comfortably. Coherence does not: its best judge exceeds
a constant-``4'' judge by ten points, and at $n=250$ the interval on $0.452$
reaches below the $0.40$ threshold. We therefore report coherence as
\emph{weakly} discriminating and note the tension this creates --- coherence
carries the paper's largest and most consistent effect while being the task on
which judges are least able to demonstrate competence. Both facts follow from
the same property, a five-point scale over gold labels that are a rounded
three-annotator mean, and we report them together rather than the convenient one
alone.

Answer-A rates confirm the swap design holds against live judges rather than only
against the artifact: $0.494$--$0.569$ on relevance and $0.500$--$0.516$ on
preference for the two judges with usable support. Sonnet's preference cell
reports $0.045$, which is not position behaviour: that cell has a $43\%$
malformed rate on $n=22$, and the figure reflects parse failure rather than a
judge that answers ``B''. We report it here rather than omitting it from the
range.

\subsection{Refusal and malformed output are separate failures, and Sonnet has both}
\label{sec:results-refusal}

Sonnet declined $24.7\%$ of relevance arm calls, reported by the provider as a
refusal with empty content. Opus refused $5.1\%$ of the same prompts and Haiku
none at all. The corpus is TREC-COVID, so the items are medical, and the effect
is a property of one judge's safety behaviour rather than of the benchmark: the
three judges saw identical prompts.

The outcome categories partition that cell's $1{,}000$ arm calls into $583$
verdicts, $249$ refusals and $168$ malformed responses. A harness that folded
refusals into malformed output would report a $41.7\%$ format-adherence failure.
The true malformed rate, computed over the arms the judge actually attempted, is
$22.4\%$ --- still a substantial and genuine instruction-following failure, and
one we report rather than let the refusal finding obscure.

Conditioning on the answered subset raises the fair objection that the three
judges are then scored on different items. We answer it by recomputing on the
common support, the 187 of 250 relevance items every judge answered under both
phrasings. The point estimates barely move: $-0.008$ for Haiku, $-0.001$ for
Opus and $+0.003$ for Sonnet. The conditioning is not what produces the
difference between these judges.

The deeper problem is not shared support. Because a rewording can itself change
whether a judge answers, refusal is an outcome of the manipulation rather than a
nuisance, and conditioning on the answered subset is selection on a
post-treatment variable. We therefore report
worst-case bounds. For this cell the conditioned JSS is identified only within
$[0.486, 0.746]$, a width of $0.26$ over $130$ unidentified pairs. That is a wide
region, and we prefer showing it to reporting a point estimate that depends on an
assumption no one can check.

\subsection{Why Sonnet's relevance cell is not evidence about wording}
\label{sec:results-sonnet-relevance}

Three facts about that cell arrive together. Its repeat ceiling is $0.789$, the
lowest we observe. Its malformed rate differs by a factor of $2.7$ between the
two arms ($0.327$ against $0.122$). And its arm-specific repeat agreement differs
by $0.209$ ($0.685$ under one template, $0.893$ under the other).

That last figure is the diagnostic we built to detect templates that are not
exchangeable, and here it fires. If one arm fails to parse three times as often
as the other, the disagreement between arms is substantially a difference in
parse surface rather than a difference in judgement, and
$\Delta\mathrm{JSS}=-0.158$ cannot be read as sensitivity to wording. The
paraphrase audit anticipated this: it warns that two of the relevance templates
state no single-token answer constraint, so their parse surfaces differ from
their partners'.

We therefore exclude this cell from the ordering and spread claims above, and
report it as what it is --- a judge whose behaviour on this corpus is dominated
by refusal and parse failure. It is included in Table~\ref{tab:main} and in the
count of intervals excluding zero, flagged, rather than deleted.

\subsection{What this evidence does not cover}
\label{sec:results-limits}

Three judges, one vendor. Every result here is within the Claude family, so any
consistency across judges is a within-family observation. The harness supports
fourteen judges across eight providers and eleven are unrun.

The preference split is complete for Haiku only. Sonnet and Opus stopped at $44$
and $170$ of $260$ rows when the API balance was exhausted mid-run; the pipeline
refuses to emit $\Delta\mathrm{JSS}$ below the declared floor of $100$ clusters,
so those cells report no endpoint rather than one computed from $17$ and $80$
clusters. Both would have been negative with intervals excluding zero, so the
exclusion costs the paper support rather than protecting it.

Intervals cluster at the item, which is the correct unit for the arms and repeats
nested inside an item --- but items are themselves nested in source records: the
$250$ relevance items come from $50$ TREC-COVID topics, the $250$ coherence items
from $92$ SummEval documents, the $250$ factuality items from $125$ TruthfulQA
questions, and the $130$ preference items from $68$ MT-Bench source records
(Section~\ref{sec:sources}). Items sharing a source share the text the
instruction wraps, so the reported intervals do not absorb source-level
correlation and are narrower than a source-clustered analysis would give. This
applies to coherence too, which carries our largest effect: its $250$ items rest
on $92$ documents, so the effective sample there is smaller than the item count
suggests. The cells nearest the threshold --- Opus on relevance and on
factuality --- are the ones this could move.

The five templates are a fixed set, not a sample. All intervals condition on
these particular strings, so the results describe sensitivity to \emph{these}
rewordings rather than to rewording in general. Relatedly, no item appears under
more than one template pair, so template identity is nested inside item identity
and no between-template comparison is identifiable.

One judge is not decoding-matched. Opus 4.7 rejects an explicit temperature
parameter and ran at the provider default while the other two ran at zero. This
is recorded on every one of its call records rather than inferred.

All measurements are at temperature zero, where a judge is closest to
deterministic and the repeat ceiling is therefore highest. Production judges are
often run above zero; there the ceiling falls, the noise term grows, and the
separation this endpoint performs becomes larger rather than smaller. We have not
measured how $\Delta\mathrm{JSS}$ behaves at $T > 0$.

All prompts are in English. Every source corpus is English-only, so whether a
judge that is stable in English remains stable under translation-equivalent
prompts is untested and is a likely place for rankings to change.

Every judge is reached through a single provider. Alias drift, replica-level
routing or an undocumented checkpoint update on that provider's side would move
every number in this paper at once, and we could not detect it by comparison
against an independent host. We record the provider-reported model string on
every call so such a change is visible after the fact, but ruling it out would
require re-running through a second host.
% ===== end 05_results =====
% ===== 06_discussion ===============================================
\section{Discussion}
\label{sec:discussion}

\subsection{What the result means for practice}

On the ordinal task, choosing a stronger judge does not buy stability. The most
and least sensitive judge on coherence differ by $0.044$
$\Delta\mathrm{JSS}$, while the tasks within a single judge span $0.162$. A
practitioner who cares about reproducibility of coherence scoring gains more by
reconsidering the response format than by upgrading the model.

We stop short of the general version of that claim. It holds decisively on
coherence and factuality, but on relevance the across-judge spread ($0.122$) is
indistinguishable from the within-judge spread it would need to beat ($0.124$),
and both figures depend on a cell we exclude as uninterpretable
(Section~\ref{sec:results-sonnet-relevance}). Establishing that task dominates
judge in general would need judges from more than one vendor.

The ordinal task is the least stable of the four, and the obvious mechanism is
the size of the answer space. A five-point scale gives a rewording more ways to
move the answer, and the weighted kappa shows most of that movement is a single
scale point.

We cannot separate that mechanism from the task, and we say so rather than
generalising. Coherence is the only ordinal task in the suite; the other three
are binary. Its position at the top of the ordering is therefore confounded with
label cardinality, and disentangling the two would require binarising the Likert
scale --- thresholding it at, say, $\geq 4$ --- and re-measuring. We have not run
that control, so ``ordinal scoring is less reproducible'' and ``coherence is less
reproducible'' are not separated by this design.

The same confound cuts the other way, and it bears directly on our factuality
null. A binary answer space quantises the judge's latent uncertainty: a judge
that moves from ``YES'' under one phrasing to a less confident ``YES'' under the
other is recorded as perfectly consistent, while the same magnitude of drift on a
five-point scale crosses a category boundary and is recorded as a flip. Binary
tasks compress the consistency signal; ordinal tasks expose it. Our factuality
null is therefore consistent with two different readings --- that judges really
are stable on factual verdicts, or that a two-valued answer space cannot show the
instability that is there --- and this design does not distinguish them. A
practitioner should not read the null as licence to trust binary judges more.

\subsection{Refusal is an evaluation failure mode, not an aside}

One judge declined a quarter of the relevance items while another, on identical
prompts, declined none. For anyone using an LLM judge over medical or otherwise
sensitive corpora this is a first-order problem: the missing verdicts are not
missing at random, and an evaluation harness that silently treats a decline as a
parse failure will report a format-adherence problem that does not exist.

It also breaks a convenience that most analyses assume. Because a rewording can
itself change whether a judge answers, refusal is an outcome of the manipulation
rather than a nuisance to be dropped, and conditioning on the answered subset is
selection on a post-treatment variable. We report worst-case bounds alongside the
conditioned estimate for this reason. Where refusals are common those bounds are
wide, and we would rather show a wide identified region than a precise number
that depends on an assumption nobody can check.

\subsection{On building a benchmark that can be wrong}

The defects we found in our own rebuild are, we think, the most transferable part
of this work. Each one passed every check that existed at the time.

A deterministic rotation assigning template pairs to items ran in lockstep with
the alternating emission of correct and incorrect answers, making the template
pair a perfect answer key --- and skewing two templates' label balance to
$75/25$, so a constant-answer judge would have appeared to have a large
template preference. Candidate responses reused across pairings carried
contradictory gold labels, so a judge reasoning correctly and consistently was
scored wrong on nearly half the preference items. A length-balancing step that
correctly removed a verbosity shortcut cut only one of the two buckets, and cut
it by vote margin, leaving the buckets matched on length and far apart on
annotator agreement.

None of these is detectable by inspecting labels, provenance, or aggregate
statistics. All three were found by asking what a heuristic with no understanding
would score, and by checking whether the artifact contained a rule that predicted
the answer. We would encourage that as routine practice: before reporting what
models score on a benchmark, report what nothing scores on it, and report it
with an interval and a support count rather than a bare number.

\subsection{Limitations}

The evidence base is three judges from one vendor family. The task ordering is
consistent across them, but that consistency is a within-family observation and
we do not claim it generalises; the harness supports fourteen judges across eight
providers and eleven are unrun. Two of the twelve judge--task cells report no
endpoint, their support having fallen below the declared floor. One judge
rejects an explicit temperature parameter and therefore ran at a provider
default while the others ran at zero; this is recorded per call rather than
inferred, and that judge's cells should be read accordingly.

The design measures sensitivity within an item across a template pair, and cannot
compare templates to one another: no item appears under more than one pair, so
template identity is nested inside item identity and the two effects are not
separable. Our equivalence checks bound specific failure modes --- label space,
polarity, construct, non-triviality --- but nothing computable from the template
strings establishes that a competent reader maps two wordings onto the same
question. Semantic drift is reduced as an alternative explanation, not
eliminated.

Finally, the effect sizes are modest on three of the four tasks, and on one they
are not distinguishable from nothing. Applying the declared $0.02$ threshold to
intervals rather than point estimates, four cells --- all three factuality cells
and Opus on relevance --- are consistent with an effect we said in advance would
not matter. We report factuality as a null. The result that survives that
standard unambiguously is the ordinal one: coherence, on all three judges, at
three to ten times the threshold.

\section{Conclusion}

We set out to measure whether LLM judges return the same verdict when asked the
same question in different words, and to do so on a benchmark that could not be
passed without judging. Across three judges and four tasks, every measurable cell
shows a real loss of agreement under meaning-preserving rewording, beyond what
the judge's own decoding variance explains. The size of that loss depends more on
the evaluation task than on the model performing it.

The benchmark, the analysis plan with its commit history, the raw judge outputs,
and the heuristic baselines that would beat the tasks are all released. So are
the defects we found in our own construction, and the measurements that exposed
them.
% ===== end 06_discussion =====

\subsubsection*{Use of AI assistance}
An AI coding assistant was used throughout this work: to implement the dataset
loaders, the metric and bootstrap code, the analysis pipeline and the test
suite; to run adversarial audits of the released files, several of which found
the construction defects reported in Sections~\ref{sec:benchmark}
and~\ref{sec:discussion}; and to draft and revise this manuscript. Every
quantitative claim was verified against committed artifacts, and the authors take
full responsibility for the content, the analysis and any remaining errors.

\subsubsection*{Reproducibility statement}
The dataset, loaders, analysis pipeline, analysis plan with its full commit
history, and every raw judge output behind the reported numbers are released at
\ifblind
\texttt{[repository URL withheld for review; the anonymised artifact accompanies
the submission]}\else\url{https://github.com/rohithreddybc/judgeSense}\fi.
Every table and every figure in Section~\ref{sec:results} is regenerated from the
committed raw outputs by a single script. Descriptive counts about the
construction, such as how many source comparisons each filter removed, are
reported from the loader's own logged output. Each split is content-hashed, and a 32-check audit gate runs
in continuous integration. Judge model identifiers, resolved decoding parameters
and the provider-reported model string are recorded on every call record, so
configuration drift is detectable after the fact rather than assumed absent.

\bibliographystyle{plainnat}          % TMLR: \bibliographystyle{tmlr}
\bibliography{references}

\end{document}
